% SPDX-License-Identifier: Apache-2.0
\section{Conclusion}
\label{sec:e-conclusion}

A companion paper merged two hardware description languages into one with
Rust-like syntax. This paper evaluated removing every construct Rust natively
provides to determine the minimal viable embedded subset.

Seventeen of twenty domain-specific constructs were eliminated, reducing the
language from a distinct grammar to four core Rust modules and procedural macros.

Two design simplifications proved especially significant. First, timeless dataflow
cannot count cycles because cycle counting requires \code{.await}, making
retiming invariants enforceable directly through Rust's type system rather than
specialized compiler checks. Second, pipelines, transaction sequences, and
concurrent processes represent a single construct (\code{async fn}) differentiated
solely by execution drivers; async state machine synthesis replaces dedicated
pipeline live-range tracking.

Two negative experimental findings guided the architecture. Arbitrary-width
arithmetic requires nightly compiler features due to const-eval limits. Furthermore,
concurrent processes cannot hold simultaneous mutable borrows of unit state;
modeling registers via interior mutability cells provides a cleaner abstraction
for shared hardware state.

To support dynamic signal selection without separate grammar parsing, the
runtime provides \code{with!}, \code{case!}, and \code{select!} macros that
lower directly into priority multiplexer trees.

While the structural features of this architecture transpose Chisel's
proven concepts into Rust with stronger static typing, the primary contribution
is unifying hardware pipelines and protocols through \code{async fn}. As
demonstrated by lowering a verified RV32I core~\cite{vreteno} at 100~MHz on
FPGA hardware, embedding TxHDL within Rust harnesses an industrial-grade type
system while concentrating design effort on novel behavioral synthesis.
